\chapter{Peripheral Arterial Bypass}
\section*{Overview}
Peripheral arterial bypass creates an alternate route around a blocked artery in a limb using a vein or synthetic graft.
\section*{Why It Is Done}
It may be required for chronic limb-threatening ischemia, severe peripheral artery disease, or selected acute arterial occlusion when endovascular treatment is unsuitable.
\section*{How It Is Performed}
The graft is connected above and below the blockage to restore blood flow to the limb. The operation may be combined with plaque removal or other revascularization.
\section*{Risks and Recovery}
Graft thrombosis, bleeding, infection, nerve injury, heart complications, kidney injury, and limb loss are possible. Antiplatelet therapy, walking, wound care, and vascular follow-up are important.
\section*{References}
\begin{enumerate}\item MedlinePlus. Peripheral Artery Disease---Surgery. U.S. National Library of Medicine; 2025.\item Current Medical Diagnosis and Treatment. McGraw-Hill; 2025.\end{enumerate}
\clearpage
